%! The Fundamental Theorem of Linear Algebra — Unit 3, Lesson 3.6 — Homework
\documentclass[10pt]{article}
\usepackage{linalg-article}
\usepackage{linalg-boxes}
\newcommand{\TallMath}[1]{$\displaystyle #1\rule[-1.4em]{0pt}{3.2em}$}
\newcommand{\vv}{\mathbf{v}}
\newcommand{\xx}{\mathbf{x}}
\newcommand{\yy}{\mathbf{y}}
\newcommand{\aaa}{\mathbf{a}}
\newcommand{\svec}{\mathbf{s}}
\newcommand{\zero}{\mathbf{0}}
\newcommand{\T}{\mathsf{T}}

\begin{document}

\pageheader{Unit 3, Lesson 3.6}{Homework}
\namedateperiod

\vspace{0.12in}

\begin{notesbox}{Practice}
The \textbf{Fundamental Theorem}: \emph{Part~1} --- $\dim C(A)=\dim C(A^{\T})=r$, $\dim N(A)=n-r$, $\dim
N(A^{\T})=m-r$; \emph{Part~2} --- $N(A)\perp C(A^{\T})$ in $\mathbb{R}^n$ and $N(A^{\T})\perp C(A)$ in
$\mathbb{R}^m$. Check a right angle with a \emph{dot product of $0$}. Always \emph{justify}.

\begin{enumerate}[leftmargin=1.9em, itemsep=11pt]
  \item \textbf{Part 1 --- dimensions.} Fill in the four dimensions for each matrix.
    \begin{center}\small
    \renewcommand{\arraystretch}{1.5}
    \begin{tabular}{@{}l c c c c@{}}
    \textbf{Matrix} & $\dim C(A)$ & $\dim C(A^{\T})$ & $\dim N(A)$ & $\dim N(A^{\T})$ \\
    \hline
    (a) $2\times5$, $r=2$ & \blank{1.1cm} & \blank{1.1cm} & \blank{1.1cm} & \blank{1.1cm} \\
    (b) $4\times4$, $r=4$ & \blank{1.1cm} & \blank{1.1cm} & \blank{1.1cm} & \blank{1.1cm} \\
    (c) $5\times3$, $r=2$ & \blank{1.1cm} & \blank{1.1cm} & \blank{1.1cm} & \blank{1.1cm} \\
    \end{tabular}
    \end{center}
  \item \textbf{Both parts, one matrix.} For $A=\begin{bmatrix} 1 & 0 & 1 \\ 2 & 1 & 3 \\ 3 & 1 & 4 \end{bmatrix}$
        (row~3 $=$ row~1 $+$ row~2), which reduces to $\begin{bmatrix} 1 & 0 & 1 \\ 0 & 1 & 1 \\ 0 & 0 & 0 \end{bmatrix}$:
        give a basis for all four subspaces (Part~1), then \emph{verify both right angles} with dot products
        (Part~2). The nullspace vector is $\svec=\begin{bsmallmatrix} -1\\-1\\1 \end{bsmallmatrix}$ and a
        left-nullspace vector is $\yy=\begin{bsmallmatrix} 1\\1\\-1 \end{bsmallmatrix}$. \vspace{2.2cm}
  \item \textbf{Read the right angle.} A matrix has row space spanned by $(1,2,2)$ and nullspace vector
        $\svec=\begin{bsmallmatrix} 2\\-2\\1 \end{bsmallmatrix}$. Show these are perpendicular, and explain in
        words what that says about the two subspaces. \writelines{2}
  \item \textbf{Explain.} Using ``$A\xx$ is computed row by row,'' explain why every vector in the nullspace
        is perpendicular to every vector in the row space. \writelines{2}
\end{enumerate}
\end{notesbox}

\begin{extensionbox}
\textbf{Orthogonal complements.}
\begin{enumerate}[label=(\alph*), leftmargin=1.8em, itemsep=5pt]
  \item \textbf{Meet only at $\zero$.} If $\vv$ is in both the row space and the nullspace, then
        $\vv\cdot\vv=0$. Explain why this forces $\vv=\zero$. (Use $\vv\cdot\vv=\|\vv\|^2$.)
  \item \textbf{The invertible case.} Let $A$ be an invertible $n\times n$ matrix. State all four dimensions
        and describe the big picture: what are the nullspace and left nullspace, and what do the row space
        and column space fill?
\end{enumerate}
\writelines{3}
\end{extensionbox}

\begin{spiralbox}
\textbf{Coming next --- Unit 4: Orthogonality.} You just discovered that a matrix's four subspaces come in two
\emph{perpendicular} pairs. Unit~4 makes that the main event: given a subspace and a point outside it, find
the closest point \emph{in} the subspace (a \emph{projection}), and use it to fit the best line to data
(\emph{least squares}). Every one of those tools rests on the right angles in today's big picture.
\end{spiralbox}

\end{document}
